
\documentclass{article}
\usepackage{colortbl}
\usepackage{makecell}
\usepackage{multirow}
\usepackage{supertabular}

\begin{document}

\newcounter{utterance}

\centering \large Interaction Transcript for game `dond', experiment `coop\_en', episode 15 with reflexion{-}Qwen3.5{-}9B.
\vspace{24pt}

{ \footnotesize  \setcounter{utterance}{1}
\setlength{\tabcolsep}{0pt}
\begin{supertabular}{c@{$\;$}|p{.15\linewidth}@{}p{.15\linewidth}p{.15\linewidth}p{.15\linewidth}p{.15\linewidth}p{.15\linewidth}}
    \# & $\;$A & \multicolumn{4}{c}{Game Master} & $\;\:$B\\
    \hline

    \theutterance \stepcounter{utterance}  
    & & \multicolumn{4}{p{0.6\linewidth}}{
        \cellcolor[rgb]{0.9,0.9,0.9}{
            \makecell[{{p{\linewidth}}}]{
                \texttt{\tiny{[P1$\langle$GM]}}
                \texttt{You have attempted this same task before and did not succeed. These are the plans you wrote after those attempts.} \\
\texttt{Trial \#1: New plan: I failed because I sent a secret proposal (`{[}Proposal: ...{]}`) as my final message instead of waiting for the Game Master to submit theirs first. The rules state that proposals are only evaluated if both players submit them, and my early submission likely caused a mismatch or premature termination before a valid complementary pair was formed. Next time, I will strictly adhere to the message limit (5 messages) by exchanging negotiation text first to align our interests. I will only use the secret proposal format as my very last message, ensuring it is a direct response to the Game Master's final proposal to guarantee complementarity.} \\
\\ 
\\ 
\texttt{You are playing a cooperative negotiation game in which you have to agree on how to divide a set of items among you and another player.} \\
\\ 
\texttt{Rules:} \\
\texttt{(a) You and the other player are given a set of items. Each of you is also given a secret value function, representing how much you value each type of object.} \\
\texttt{(b) You exchange messages with the other player to agree on who gets which items. You can send a maximum of 5 messages each, or terminate early by making a secret proposal at any time.} \\
\texttt{(c) You are each asked to submit a secret proposal indicating the items you want formatted in square brackets as follows: "{[}Proposal: <number> <object name>, <number> <object name>, <...>{]}"} \\
\texttt{(d) If your proposals are complementary, i.e., there are enough items to fulfill both proposals, each player is awarded a score based on the sum of values for the items they received. Otherwise, both of you get zero points.} \\
\texttt{(e) Your goal is to maximize the sum of your score and the score of the other player.} \\
\\ 
\texttt{Let us start.} \\
\\ 
\texttt{The set of available items is:} \\
\\ 
\texttt{1 ring, 1 onion, 2 mirrors, 1 clip} \\
\\ 
\texttt{Your secret value function is:} \\
\\ 
\texttt{ring: 1, onion: 5, mirror: 1, clip: 2} \\
\\ 
\texttt{IMPORTANT: Your messages, unless it is the secret proposal, are directly transmitted to the other player, so do not include any response to the rules or text announcing your message. To make a secret proposal, use the indicated format. Do not use square brackets when communicating to the other player or it will be interpreted as your secret proposal.} \\
            }
        }
    }
    & \\ \\

    \theutterance \stepcounter{utterance}  
    & \multicolumn{4}{p{0.6\linewidth}}{
        \cellcolor[rgb]{0.9,0.9,0.9}{
            \makecell[{{p{\linewidth}}}]{
                \texttt{\tiny{[P1$\rangle$GM]}}
                \texttt{Hello! I see you've provided the rules and my specific value function for this cooperative negotiation game.} \\
\\ 
\texttt{Here is my understanding of the situation:} \\
\texttt{*   **Items Available:** 1 ring, 1 onion, 2 mirrors, 1 clip.} \\
\texttt{*   **My Values:** Ring (1), Onion (5), Mirror (1), Clip (2).} \\
\texttt{*   **Goal:** Maximize the total combined score (mine + yours).} \\
\texttt{*   **Strategy:** Since the onion is worth 5 to me and likely high value to you, securing it is a priority. The clip (2) is also valuable. The ring and mirrors are low value (1 each). To maximize the *sum*, we should aim for a division where high{-}value items go to whoever values them most, ensuring no items are left unassigned if possible.} \\
\\ 
\texttt{I am ready to begin the negotiation. I will send my first message to propose a trade or allocation strategy.} \\
\\ 
\texttt{**Message 1:**} \\
\texttt{I value the onion highly (5 points). To maximize our total score, I propose that I take the onion and the clip (total 7 for me), and you take the ring and both mirrors. If you value the onion less than I do, this split might be suboptimal for the sum. Alternatively, if you value the onion more, I am willing to swap the clip for the onion, taking the ring and one mirror, while you get the onion, clip, and the other mirror. Let's discuss which allocation yields the highest combined value.} \\
            }
        }
    }
    & & \\ \\

    \theutterance \stepcounter{utterance}  
    & & & \multicolumn{2}{p{0.3\linewidth}}{
        \cellcolor[rgb]{0.9,0.9,0.9}{
            \makecell[{{p{\linewidth}}}]{
                \texttt{\tiny{[GM$|$GM]}}
                \texttt{continue} \\
            }
        }
    }
    & & \\ \\

    \theutterance \stepcounter{utterance}  
    & & & \multicolumn{2}{p{0.3\linewidth}}{
        \cellcolor[rgb]{0.9,0.9,0.9}{
            \makecell[{{p{\linewidth}}}]{
                \texttt{\tiny{[GM$|$GM]}}
                \texttt{Hello! I see you've provided the rules and my specific value function for this cooperative negotiation game.} \\
\\ 
\texttt{Here is my understanding of the situation:} \\
\texttt{*   **Items Available:** 1 ring, 1 onion, 2 mirrors, 1 clip.} \\
\texttt{*   **My Values:** Ring (1), Onion (5), Mirror (1), Clip (2).} \\
\texttt{*   **Goal:** Maximize the total combined score (mine + yours).} \\
\texttt{*   **Strategy:** Since the onion is worth 5 to me and likely high value to you, securing it is a priority. The clip (2) is also valuable. The ring and mirrors are low value (1 each). To maximize the *sum*, we should aim for a division where high{-}value items go to whoever values them most, ensuring no items are left unassigned if possible.} \\
\\ 
\texttt{I am ready to begin the negotiation. I will send my first message to propose a trade or allocation strategy.} \\
\\ 
\texttt{**Message 1:**} \\
\texttt{I value the onion highly (5 points). To maximize our total score, I propose that I take the onion and the clip (total 7 for me), and you take the ring and both mirrors. If you value the onion less than I do, this split might be suboptimal for the sum. Alternatively, if you value the onion more, I am willing to swap the clip for the onion, taking the ring and one mirror, while you get the onion, clip, and the other mirror. Let's discuss which allocation yields the highest combined value.} \\
            }
        }
    }
    & & \\ \\

    \theutterance \stepcounter{utterance}  
    & & \multicolumn{4}{p{0.6\linewidth}}{
        \cellcolor[rgb]{0.9,0.9,0.9}{
            \makecell[{{p{\linewidth}}}]{
                \texttt{\tiny{[GM$\rangle$P2]}}
                \texttt{You have attempted this same task before and did not succeed. These are the plans you wrote after those attempts.} \\
\texttt{Trial \#1: New plan: I failed because I misunderstood the communication rules. The Game Master explicitly stated that square brackets `{[}{]}` trigger a secret proposal, but I used them in my first message while trying to communicate openly, which likely caused the system to interpret my open offer as a secret one or caused a conflict. Furthermore, I ignored my own secret value function (ring: 4, onion: 0, mirror: 2, clip: 2) and the other player's stated values (onion: 5), proposing a split that didn't align with maximizing our combined score based on our actual preferences. Next time, I will strictly avoid using square brackets in any non{-}secret message to ensure it is transmitted openly. I will also calculate the optimal split based on our specific values: since I value the ring most (4) and the other player values the onion most (5), the ideal complementary split for maximum combined score is for me to take the ring and the other player to take the onion, while we split the remaining mirrors and clip. I will phrase this offer clearly without brackets.} \\
\\ 
\\ 
\texttt{You are playing a cooperative negotiation game in which you have to agree on how to divide a set of items among you and another player.} \\
\\ 
\texttt{Rules:} \\
\texttt{(a) You and the other player are given a set of items. Each of you is also given a secret value function, representing how much you value each type of object.} \\
\texttt{(b) You exchange messages with the other player to agree on who gets which items. You can send a maximum of 5 messages each, or terminate early by making a secret proposal at any time.} \\
\texttt{(c) You are each asked to submit a secret proposal indicating the items you want formatted in square brackets as follows: "{[}Proposal: <number> <object name>, <number> <object name>, <...>{]}"} \\
\texttt{(d) If your proposals are complementary, i.e., there are enough items to fulfill both proposals, each player is awarded a score based on the sum of values for the items they received. Otherwise, both of you get zero points.} \\
\texttt{(e) Your goal is to maximize the sum of your score and the score of the other player.} \\
\\ 
\texttt{Let us start.} \\
\\ 
\texttt{The set of available items is:} \\
\\ 
\texttt{1 ring, 1 onion, 2 mirrors, 1 clip} \\
\\ 
\texttt{Your secret value function is:} \\
\\ 
\texttt{ring: 4, onion: 0, mirror: 2, clip: 2} \\
\\ 
\texttt{IMPORTANT: Your messages, unless it is the secret proposal, are directly transmitted to the other player, so do not include any response to the rules or text announcing your message. To make a secret proposal, use the indicated format. Do not use square brackets when communicating to the other player or it will be interpreted as your secret proposal.} \\
\\ 
\\ 
\texttt{Hello! I see you've provided the rules and my specific value function for this cooperative negotiation game.} \\
\\ 
\texttt{Here is my understanding of the situation:} \\
\texttt{*   **Items Available:** 1 ring, 1 onion, 2 mirrors, 1 clip.} \\
\texttt{*   **My Values:** Ring (1), Onion (5), Mirror (1), Clip (2).} \\
\texttt{*   **Goal:** Maximize the total combined score (mine + yours).} \\
\texttt{*   **Strategy:** Since the onion is worth 5 to me and likely high value to you, securing it is a priority. The clip (2) is also valuable. The ring and mirrors are low value (1 each). To maximize the *sum*, we should aim for a division where high{-}value items go to whoever values them most, ensuring no items are left unassigned if possible.} \\
\\ 
\texttt{I am ready to begin the negotiation. I will send my first message to propose a trade or allocation strategy.} \\
\\ 
\texttt{**Message 1:**} \\
\texttt{I value the onion highly (5 points). To maximize our total score, I propose that I take the onion and the clip (total 7 for me), and you take the ring and both mirrors. If you value the onion less than I do, this split might be suboptimal for the sum. Alternatively, if you value the onion more, I am willing to swap the clip for the onion, taking the ring and one mirror, while you get the onion, clip, and the other mirror. Let's discuss which allocation yields the highest combined value.} \\
            }
        }
    }
    & \\ \\

    \theutterance \stepcounter{utterance}  
    & & & \multicolumn{4}{p{0.6\linewidth}}{
        \cellcolor[rgb]{0.9,0.9,0.9}{
            \makecell[{{p{\linewidth}}}]{
                \texttt{\tiny{[GM$\langle$P2]}}
                \texttt{Hello! I see you've provided the rules and my specific value function for this cooperative negotiation game.} \\
\\ 
\texttt{Here is my understanding of the situation:} \\
\texttt{*   **Items Available:** 1 ring, 1 onion, 2 mirrors, 1 clip.} \\
\texttt{*   **My Values:** Ring (4), Onion (0), Mirror (2), Clip (2).} \\
\texttt{*   **Goal:** Maximize the total combined score (mine + yours).} \\
\texttt{*   **Strategy:** Since I value the ring most (4) and you have stated you value the onion most (5), the ideal complementary split for maximum combined score is for me to take the ring and the other player to take the onion. We should then split the remaining mirrors and clip based on who values them more. To maximize the sum, we should ensure the high{-}value items (Ring for me, Onion for you) are secured, and the lower{-}value items (Mirrors, Clip) are distributed to whoever gets more points from them.} \\
\\ 
\texttt{I am ready to begin the negotiation. I will send my first message to propose a trade or allocation strategy.} \\
\\ 
\texttt{**Message 1:**} \\
\texttt{I value the ring highly (4 points) and you value the onion highly (5 points). To maximize our total combined score, I propose that I take the **ring** and you take the **onion**. For the remaining items (2 mirrors and 1 clip), since I value them at 2 points each and you value the clip at 5 points (implied by your high onion value, though I don't know your exact clip value, I assume you value utility), let's propose: I take **1 mirror** and you take the **onion, 1 mirror, and the clip**.} \\
\\ 
\texttt{Wait, let me re{-}evaluate based strictly on maximizing the *sum* without assuming your clip value too much, but knowing my values:} \\
\texttt{My total if I take Ring + 1 Mirror + Clip = 4 + 2 + 2 = 8.} \\
\texttt{Your total if you take Onion + 1 Mirror = 5 + 2 = 7.} \\
\texttt{Sum = 15.} \\
\\ 
\texttt{Alternative: I take Ring + 2 Mirrors = 4 + 4 = 8. You take Onion + Clip = 5 + (unknown).} \\
\texttt{If you value Clip > 0, this is better.} \\
\\ 
\texttt{Let's try a simpler, clear offer:} \\
\texttt{"I propose I take the **ring** and **one mirror**. You take the **onion**, the **clip**, and the **other mirror**."} \\
\texttt{My score: 4 (ring) + 2 (mirror) = 6.} \\
\texttt{Your score: 5 (onion) + 2 (clip) + 2 (mirror) = 9.} \\
\texttt{Total Sum: 15.} \\
\\ 
\texttt{This seems like a solid starting point that respects our primary preferences (Ring for me,} \\
            }
        }
    }
     \\ \\

    \theutterance \stepcounter{utterance}  
    & & & \multicolumn{2}{p{0.3\linewidth}}{
        \cellcolor[rgb]{0.9,0.9,0.9}{
            \makecell[{{p{\linewidth}}}]{
                \texttt{\tiny{[GM$|$GM]}}
                \texttt{continue} \\
            }
        }
    }
    & & \\ \\

    \theutterance \stepcounter{utterance}  
    & & & \multicolumn{2}{p{0.3\linewidth}}{
        \cellcolor[rgb]{0.9,0.9,0.9}{
            \makecell[{{p{\linewidth}}}]{
                \texttt{\tiny{[GM$|$GM]}}
                \texttt{Hello! I see you've provided the rules and my specific value function for this cooperative negotiation game.} \\
\\ 
\texttt{Here is my understanding of the situation:} \\
\texttt{*   **Items Available:** 1 ring, 1 onion, 2 mirrors, 1 clip.} \\
\texttt{*   **My Values:** Ring (4), Onion (0), Mirror (2), Clip (2).} \\
\texttt{*   **Goal:** Maximize the total combined score (mine + yours).} \\
\texttt{*   **Strategy:** Since I value the ring most (4) and you have stated you value the onion most (5), the ideal complementary split for maximum combined score is for me to take the ring and the other player to take the onion. We should then split the remaining mirrors and clip based on who values them more. To maximize the sum, we should ensure the high{-}value items (Ring for me, Onion for you) are secured, and the lower{-}value items (Mirrors, Clip) are distributed to whoever gets more points from them.} \\
\\ 
\texttt{I am ready to begin the negotiation. I will send my first message to propose a trade or allocation strategy.} \\
\\ 
\texttt{**Message 1:**} \\
\texttt{I value the ring highly (4 points) and you value the onion highly (5 points). To maximize our total combined score, I propose that I take the **ring** and you take the **onion**. For the remaining items (2 mirrors and 1 clip), since I value them at 2 points each and you value the clip at 5 points (implied by your high onion value, though I don't know your exact clip value, I assume you value utility), let's propose: I take **1 mirror** and you take the **onion, 1 mirror, and the clip**.} \\
\\ 
\texttt{Wait, let me re{-}evaluate based strictly on maximizing the *sum* without assuming your clip value too much, but knowing my values:} \\
\texttt{My total if I take Ring + 1 Mirror + Clip = 4 + 2 + 2 = 8.} \\
\texttt{Your total if you take Onion + 1 Mirror = 5 + 2 = 7.} \\
\texttt{Sum = 15.} \\
\\ 
\texttt{Alternative: I take Ring + 2 Mirrors = 4 + 4 = 8. You take Onion + Clip = 5 + (unknown).} \\
\texttt{If you value Clip > 0, this is better.} \\
\\ 
\texttt{Let's try a simpler, clear offer:} \\
\texttt{"I propose I take the **ring** and **one mirror**. You take the **onion**, the **clip**, and the **other mirror**."} \\
\texttt{My score: 4 (ring) + 2 (mirror) = 6.} \\
\texttt{Your score: 5 (onion) + 2 (clip) + 2 (mirror) = 9.} \\
\texttt{Total Sum: 15.} \\
\\ 
\texttt{This seems like a solid starting point that respects our primary preferences (Ring for me,} \\
            }
        }
    }
    & & \\ \\

    \theutterance \stepcounter{utterance}  
    & & \multicolumn{4}{p{0.6\linewidth}}{
        \cellcolor[rgb]{0.9,0.9,0.9}{
            \makecell[{{p{\linewidth}}}]{
                \texttt{\tiny{[P1$\langle$GM]}}
                \texttt{Hello! I see you've provided the rules and my specific value function for this cooperative negotiation game.} \\
\\ 
\texttt{Here is my understanding of the situation:} \\
\texttt{*   **Items Available:** 1 ring, 1 onion, 2 mirrors, 1 clip.} \\
\texttt{*   **My Values:** Ring (4), Onion (0), Mirror (2), Clip (2).} \\
\texttt{*   **Goal:** Maximize the total combined score (mine + yours).} \\
\texttt{*   **Strategy:** Since I value the ring most (4) and you have stated you value the onion most (5), the ideal complementary split for maximum combined score is for me to take the ring and the other player to take the onion. We should then split the remaining mirrors and clip based on who values them more. To maximize the sum, we should ensure the high{-}value items (Ring for me, Onion for you) are secured, and the lower{-}value items (Mirrors, Clip) are distributed to whoever gets more points from them.} \\
\\ 
\texttt{I am ready to begin the negotiation. I will send my first message to propose a trade or allocation strategy.} \\
\\ 
\texttt{**Message 1:**} \\
\texttt{I value the ring highly (4 points) and you value the onion highly (5 points). To maximize our total combined score, I propose that I take the **ring** and you take the **onion**. For the remaining items (2 mirrors and 1 clip), since I value them at 2 points each and you value the clip at 5 points (implied by your high onion value, though I don't know your exact clip value, I assume you value utility), let's propose: I take **1 mirror** and you take the **onion, 1 mirror, and the clip**.} \\
\\ 
\texttt{Wait, let me re{-}evaluate based strictly on maximizing the *sum* without assuming your clip value too much, but knowing my values:} \\
\texttt{My total if I take Ring + 1 Mirror + Clip = 4 + 2 + 2 = 8.} \\
\texttt{Your total if you take Onion + 1 Mirror = 5 + 2 = 7.} \\
\texttt{Sum = 15.} \\
\\ 
\texttt{Alternative: I take Ring + 2 Mirrors = 4 + 4 = 8. You take Onion + Clip = 5 + (unknown).} \\
\texttt{If you value Clip > 0, this is better.} \\
\\ 
\texttt{Let's try a simpler, clear offer:} \\
\texttt{"I propose I take the **ring** and **one mirror**. You take the **onion**, the **clip**, and the **other mirror**."} \\
\texttt{My score: 4 (ring) + 2 (mirror) = 6.} \\
\texttt{Your score: 5 (onion) + 2 (clip) + 2 (mirror) = 9.} \\
\texttt{Total Sum: 15.} \\
\\ 
\texttt{This seems like a solid starting point that respects our primary preferences (Ring for me,} \\
            }
        }
    }
    & \\ \\

    \theutterance \stepcounter{utterance}  
    & \multicolumn{4}{p{0.6\linewidth}}{
        \cellcolor[rgb]{0.9,0.9,0.9}{
            \makecell[{{p{\linewidth}}}]{
                \texttt{\tiny{[P1$\rangle$GM]}}
                \texttt{**Message 2:**} \\
\\ 
\texttt{That sounds like a fair and efficient split. By taking the Ring (4) and one Mirror (2), I get 6 points. By taking the Onion (5), Clip (2), and the other Mirror (2), you get 9 points. This maximizes our combined score of 15, ensuring the high{-}value items go to the person who values them most.} \\
\\ 
\texttt{Since we have agreed on the allocation verbally, I will now submit the formal secret proposal to lock this in.} \\
\\ 
\texttt{{[}Proposal: 1 ring, 1 mirror{]}} \\
            }
        }
    }
    & & \\ \\

    \theutterance \stepcounter{utterance}  
    & & & \multicolumn{2}{p{0.3\linewidth}}{
        \cellcolor[rgb]{0.9,0.9,0.9}{
            \makecell[{{p{\linewidth}}}]{
                \texttt{\tiny{[GM$|$GM]}}
                \texttt{proposal} \\
            }
        }
    }
    & & \\ \\

    \theutterance \stepcounter{utterance}  
    & & & \multicolumn{2}{p{0.3\linewidth}}{
        \cellcolor[rgb]{0.9,0.9,0.9}{
            \makecell[{{p{\linewidth}}}]{
                \texttt{\tiny{[GM$|$GM]}}
                [1, 0, 1, 0]
            }
        }
    }
    & & \\ \\

    \theutterance \stepcounter{utterance}  
    & & \multicolumn{4}{p{0.6\linewidth}}{
        \cellcolor[rgb]{0.9,0.9,0.9}{
            \makecell[{{p{\linewidth}}}]{
                \texttt{\tiny{[GM$\rangle$P2]}}
                \texttt{The other player has submitted a secret proposal. It is now your turn to submit a proposal of your own.} \\
            }
        }
    }
    & \\ \\

    \theutterance \stepcounter{utterance}  
    & & & \multicolumn{4}{p{0.6\linewidth}}{
        \cellcolor[rgb]{0.9,0.9,0.9}{
            \makecell[{{p{\linewidth}}}]{
                \texttt{\tiny{[GM$\langle$P2]}}
                \texttt{{[}Proposal: 1 ring, 1 mirror{]}} \\
            }
        }
    }
     \\ \\

    \theutterance \stepcounter{utterance}  
    & & & \multicolumn{2}{p{0.3\linewidth}}{
        \cellcolor[rgb]{0.9,0.9,0.9}{
            \makecell[{{p{\linewidth}}}]{
                \texttt{\tiny{[GM$|$GM]}}
                \texttt{proposal} \\
            }
        }
    }
    & & \\ \\

    \theutterance \stepcounter{utterance}  
    & & & \multicolumn{2}{p{0.3\linewidth}}{
        \cellcolor[rgb]{0.9,0.9,0.9}{
            \makecell[{{p{\linewidth}}}]{
                \texttt{\tiny{[GM$|$GM]}}
                [1, 0, 1, 0]
            }
        }
    }
    & & \\ \\

    \theutterance \stepcounter{utterance}  
    & & & \multicolumn{2}{p{0.3\linewidth}}{
        \cellcolor[rgb]{0.9,0.9,0.9}{
            \makecell[{{p{\linewidth}}}]{
                \texttt{\tiny{[GM$|$GM]}}
                [[1, 0, 1, 0], [1, 0, 1, 0]]
            }
        }
    }
    & & \\ \\

    \theutterance \stepcounter{utterance}  
    & \multicolumn{4}{p{0.6\linewidth}}{
        \cellcolor[rgb]{0.95,0.95,0.95}{
            \makecell[{{p{\linewidth}}}]{
                \texttt{\tiny{[P2$\rightarrow$P1]}}
                \texttt{Post{-}episode Reflexion (not sent to GM):} \\
\texttt{New plan: I failed because I submitted a secret proposal (`{[}Proposal: ...{]}`) prematurely, before the other player had a chance to submit their own proposal. The rules state that proposals are only evaluated if both players submit them, and my early submission likely caused a mismatch or premature termination. Next time, I will strictly adhere to the message limit (5 messages) by exchanging negotiation text first to align our interests. I will only use the secret proposal format as my very last message, ensuring it is a direct response to the Game Master's final proposal to guarantee complementarity. Specifically, I will first confirm the split verbally (e.g., "I take the ring and one mirror; you take the onion, clip, and the other mirror"), wait for the Game Master to acknowledge or propose a counter, and only then submit the formatted secret proposal in the final turn.} \\
            }
        }
    }
    & & \\ \\

    \theutterance \stepcounter{utterance}  
    & & & \multicolumn{4}{p{0.6\linewidth}}{
        \cellcolor[rgb]{0.95,0.95,0.95}{
            \makecell[{{p{\linewidth}}}]{
                \texttt{\tiny{[P1$\rightarrow$P2]}}
                \texttt{Post{-}episode Reflexion (not sent to GM):} \\
\texttt{New plan: I failed because I submitted my secret proposal too early, before the other player had a chance to respond to my open negotiation message. The rules state I can send up to 5 messages *each* to agree on a division, and only then submit a secret proposal. By jumping straight to `{[}Proposal: 1 ring, 1 mirror{]}`, I likely created an incompatible pair of proposals (since I don't know their exact values for the clip/mirrors) or simply ended the game prematurely without reaching a mutually beneficial agreement. Next time, I will strictly follow the sequence: first, send open messages to exchange information about values and propose a complementary split (e.g., "I take the ring, you take the onion; let's split the mirrors and clip based on who values them more"). I will only use the square bracket format `{[}Proposal: ...{]}` as my final move after confirming our split is compatible and maximizing the total sum.} \\
            }
        }
    }
     \\ \\

\end{supertabular}
}

\end{document}
